\chapter{Cold Hands and Feet (Raynaud Phenomenon)}

\section*{Definition}
A vasospastic disorder characterized by episodic digital ischemia triggered by cold or stress, causing well-demarcated color changes (white, blue, red) in the fingers and/or toes, often with numbness, tingling, and pain during rewarming.

\section*{Pathophysiology}
Raynaud phenomenon involves exaggerated vasoconstriction of digital arteries and arterioles in response to cold or emotional stress. Primary Raynaud is idiopathic, benign, and associated with hyperactivation of alpha-2 adrenergic receptors on digital vessels. Secondary Raynaud involves underlying pathology (scleroderma, lupus) causing structural vessel damage, intimal hyperplasia, and obliterative microvascular disease, leading to more severe ischemia and potential digital ulceration.

\section*{Etiology}
\begin{itemize}
  \item Primary Raynaud disease (benign, no underlying cause, onset before age 30)
  \item Connective tissue diseases (scleroderma, systemic lupus, mixed connective tissue disease)
  \item Cryoglobulinemia (abnormal cold-precipitating proteins)
  \item Buerger disease (thromboangiitis obliterans, linked to smoking)
  \item Atherosclerosis or peripheral artery disease
  \item Thoracic outlet syndrome (neurovascular compression)
  \item Medications (beta-blockers, ergotamines, chemotherapy bleomycin, cisplatin)
  \item Repetitive vibration injury (hand-arm vibration syndrome from power tools)
  \item Hypothyroidism (reduced metabolic rate, cold intolerance)
  \item Carpal tunnel syndrome (associated with vasomotor symptoms)
\end{itemize}

\section*{Indications for Evaluation}
Seek care if:
\begin{itemize}
  \item Severe or worsening symptoms
  \item Unilateral symptoms, skin ulceration or gangrene, symptoms beginning after age 30, or associated with rash, joint pain, or systemic symptoms
  \item Accompanied by other concerning signs
\end{itemize}

\section*{Related Symptoms}
\begin{itemize}
  \item Numbness
  \item Tingling
  \item Finger swelling
  \item Joint pain
  \item Digital ulcers
\end{itemize}
\textbf{Note:} Always consider serious underlying conditions when evaluating persistent cold hands and feet (raynaud phenomenon).

\section*{Self-Care / Home Management}
\begin{itemize}
  \item Rest and avoid activities that may aggravate the condition.
  \item Maintain adequate hydration and a balanced diet.
  \item Over-the-counter remedies may provide symptomatic relief (consult a pharmacist or doctor).
  \item Monitor symptoms and seek professional advice if they persist or worsen.
  \item Avoid self-medicating with prescription medications without medical guidance.
\end{itemize}

\section*{Diagnostic Approach}
\begin{itemize}
  \item A thorough history and physical examination are the first steps in evaluation.
  \item Laboratory tests (blood work, urinalysis) may be ordered based on clinical suspicion.
  \item Imaging studies (X-ray, ultrasound, CT, MRI) may be indicated when structural pathology is suspected.
  \item Specialist referral is arranged when the diagnosis remains uncertain or requires specific expertise.
\end{itemize}

\section*{Treatment / Management Overview}
\begin{itemize}
  \item Treatment targets the underlying cause identified through diagnostic evaluation.
  \item Supportive care and symptom management are provided while awaiting definitive therapy.
  \item Medications, lifestyle modifications, and physical therapies may be prescribed as appropriate.
  \item Follow-up ensures treatment effectiveness and allows timely adjustments to the care plan.
\end{itemize}

\clearpage